\section{PPU}

% Architektura
%   - struktura PPU
%   - presna reprezentace vsech registru a pameti
%   - vyuziti Loopyho modelu pro scrolling
% IO/komunikace
%   - ukazka ppu_write() a ppu_read() funkci
%   - implementace OAM DMA
% Rendrovani
%   - ukazka ppu_run_until() funkce
%     - pocitani cyklu
%   - buferizovani pixelu
%   - algoritmus fetchovani tilu (background a sprity)
%     - LUT tabulky
%   - algoritmus rendrovani backgroundu
%     - realizace scrollingu
%   - algoritmus rendrovani spritu
%   - algoritmus scanu spritu

Funkce \texttt{ppu\_run\_until} (viz výpis \ref{lst:ppu_run_until} aktualizuje stav PPU a posouvá ho
po cyklové ose dokud nedožene stav CPU (parametr \texttt{target\_cycle}).

\begin{listing}[htbp]
\begin{minted}{c}
void ppu_run_until(Ppu *self, CycleCounter target_cycle) {
    while (self->cycles < target_cycle) {
        unsigned cycles_rem = SCANLINE_END - self->dot;
        unsigned cycles_to_run = target_cycle - self->cycles;
        unsigned step =
            (cycles_to_run < cycles_rem) ? cycles_to_run : cycles_rem;
        unsigned prev_dot = self->dot;
        unsigned next_dot = self->dot + step;

        /* Aktualizace stavu PPU... */

        self->cycles += step;
        self->dot = next_dot;
        /* Logika přechodu mezi scanline */
        if (self->dot >= 341) {
            /* inkrementace scanline, případné nastavení VBlank, ... */
        }
    }
}
\end{minted}
\caption{Funkce pro aktualizaci stavu PPU, implementovaná pomocí catch-up mechanismu.}
\label{lst:ppu_run_until}
\end{listing}

Namísto inkrementální emulace po jednotlivých tečkách využívá funkce \texttt{ppu\_run\_until}
principu skoků k významným událostem v rámci řádku. Algoritmus vypočítá vzdálenost k nejbližšímu
bodu zájmu (např. cyklus odpovídající cyklu CPU nebo konec řádku) a provede aktualizaci stavu v
blocích. Tento přístup výrazně šetří instrukční cykly procesoru RP2350.

\label{ppu_scanline_callback}
Když PPU dokončí vykreslování řádku, vyvolá callback, kde cílový kód může přímo poslat řádek na
displej nebo provést další operace. To dovoluje PPU být nezávislým na způsobu vykreslování na cílové
architektuře. Pixely se vykreslují do vnitřního bufferu.

\subsection{Vykreslování pozadí}

Proces generování obrazu v emulátoru musí být velice efektivní, aby RP2350 dokázal udržet stabilní
snímkovou frekvenci. Funkce \texttt{render\_bg} (výpis \ref{lst:render_bg}) realizuje vykreslování
dlaždic pozadí s pixelovou granulací. Ačkoliv se data z paměti čtou po řádcích dlaždic (8 pixelů),
samotný zápis do bufferu \texttt{scanline\_buf} může probíhat po blocích. Tento přístup je také
nezbytný pro správnou realizaci jemného horizontálního posunu (\textit{fine X scrolling}).

\begin{listing}[htbp]
\begin{minted}{c}
void render_bg(Ppu *self, unsigned start, unsigned end) {
    unsigned curr = start;
    /* Výpočet počátečního offsetu v rámci první dlaždice */
    unsigned offset = (curr + self->x) & 7;
    uint8_t *dest = &self->scanline_buf[curr];
    while (curr < end) {
        /* Výpočet úseku, který zbývá vykreslit v aktuální dlaždici */
        unsigned rem = 8 - offset;
        unsigned count = (rem > (end - curr)) ? (end - curr) : rem;
        /* Načtení dat dlaždice a atributů z pamětí */
        TileRow tile = fetch_tile(self);
        /* Bitový posun pro zarovnání pixelů podle jemného X posunu */
        tile >>= (offset << 2);
        /* Vnitřní smyčka pro zápis pixelů */
        for (unsigned i = 0; i < count; ++i) {
            *dest++ = tile & 0x0f; /* zápis 4bitové barvy */
            tile >>= 4;
            ++curr;
        }
        /* Další dlaždice již začínají od nultého pixelu */
        offset = 0;
        /* Kontrola přechodu na další dlaždici (hrubý X posun) */
        if (!((curr + self->x) & 7)) {
            increment_coarse_x(self);
        }
    }
}
\end{minted}
\caption{Funkce pro vykreslování dlaždicových řádku pozadí.}
\label{lst:render_bg}
\end{listing}

V rámci implementace byly použity tři klíčové optimalizační techniky:

\begin{itemize}
    \item Logika jemného posunu je integrována přímo do kódu vykreslování pomocí bitových posunů.
        Tím odpadá nutnost dodatečného ořezávání a posunu obrazu.
    \item V kritické smyčce \texttt{for}, která zapisuje pixely do \texttt{scanline\_buf}, se využívá
        pointerová aritmetika (\texttt{*dest++}) namísto přímého indexování pole pomocí \texttt{curr}.
        Slouží to jako optimalizační nápověda pro kompilátor, aby efektivněji využil registry a
        instrukce procesoru RP2350. Tím odpadává nutnost opětovného výpočtu adresy v každé iteraci
        smyčky.
    \item „Balení“ řádku dlaždice ve 32bitovém typu: tradiční reprezentace barev pomocí RGB v NES je
        pro moderní procesory nevýhodná. Byla využita technika balení osmi pixelů do jedné 32bitové
        proměnné (typ \texttt{TileRow}). Každý 4bitový nibble v této proměnné nese kompletní
        informaci o indexu palety i barvy daného pixelu. To umožňuje provádět jemný horizontální
        posun pomocí jediného bitového posunu nad celým registrem, namísto osmi samostatných
        operací. Funkce \texttt{fetch\_tile()}, která balí jeden řádek dlaždice do 32bitového typu,
        bude popsáná v další sekcí.
\end{itemize}

\subsection{Příprava dat pro vykreslování pozadí}

Funkce \texttt{fetch\_tile()} (výpis \ref{lst:fetch_tile}) slouží k načtení celé dlaždice z VRAM
podle adresy, nastavené interním 16bitovým registrem \texttt{v}.

\begin{listing}[htbp]
\begin{minted}{c}
TileRow fetch_tile(const Ppu *self) {
    uint16_t v = self->v; /* Aktuální adresní registr */
    /* 1. Výpočet adresy a načtení indexu dlaždice z Name Table */
    uint16_t addr = PPU_NT0 | (v & (VRAM_COARSE_X | VRAM_COARSE_Y | VRAM_NT_SEL));
    uint8_t tile_idx = self->nt[(addr >> 10) & 3][addr & 0x3ff];
    /* 2. Načtení atributů pro danou oblast */
    uint16_t attr_addr = calc_attr_addr(v);
    uint8_t attr = self->nt[(attr_addr >> 10) & 3][attr_addr & 0x3ff];
    /* 3. Výpočet posunu pro extrakci 2 bitů palety z bajtu atributů */
    unsigned shift = ((v >> 4) & 4) | (v & 2);
    uint8_t palette_idx = (attr >> shift) & 3;
    /* 4. Adresace Pattern Table a načtení dvou bitových rovin */
    uint16_t row_addr = get_pattern_table_bg(self) +
        (tile_idx << 4) + vram_get_fine_y(v);
    uint8_t row_low = self->chr[row_addr];
    uint8_t row_high = self->chr[row_addr + 8];
    /* 5. Syntéza 32bitového slova pomocí předpočítaných tabulek (LUT).
     * Transformace bitových rovin na rozprostřenou masku pixelů */
    uint32_t raw_pixels = row_high_lut[row_high] | row_low_lut[row_low];
    /* Aplikace indexu palety na všechny pixely v bloku */
    uint32_t attr_mask = attr_mask_lut[palette_idx];
    /* 6. Zabalení řádku do 32bitového typu TileRow (8 pixelů po 4 bitech) */
    return raw_pixels | attr_mask;
}
\end{minted}
\caption{Extrakce řádku dlaždice a následná jeho syntézu do 32bitového typu \texttt{TileRow} pomocí
LUT.}
\label{lst:fetch_tile}
\end{listing}

Jak je patrné, využití LUT tabulek v této funkcí slouží k úplnému zamezení redundantním výpočtům v
kritické cestě emulátoru. Namísto procesorově náročné manipulace s jednotlivými pixely je
transformace grafických dat redukována na několik operací čtení z paměti a bitových operací.

\subsection{Vykreslování spritů}

Po vykreslení bloku pixelů pozadí, proběhne algoritmus vykreslování spritů, který pracuje s
jednotlivými sprity na řádku (pokud jsou, nebo vykreslování spritů je zapnuto). Funkce
\texttt{render\_sprites} (viz výpis \ref{lst:render_sprites}) iteruje frontu spritů, řeší Sprite 0
Hit a prioritu vůči pozadí. Výsledné pixely zapisuje do \texttt{scanline\_buf}.

\begin{listing}[htbp]
\begin{minted}{c}
for (unsigned x = draw_start; x < draw_end; ++x) {
    uint8_t px = tile & 0x0f; /* Pixel spritu */
    uint8_t bg_px = *dest & 0x0f; /* Existující pixel pozadí v bufferu */
    bool bg_opaque = (bg_px & 0x03) != 0;
    /* 2. Je pixel spritu neprůhledný? */
    if ((px & 0x03) != 0) {
        /* 3. Detekce kolize pro Sprite 0 Hit */
        if (spr->is_0 && bg_opaque && x < 255) {
            self->stat |= PPU_STAT_ZERO_HIT;
        }
        /* 4. Zápis pixelu při splnění prioritních podmínek */
        if (!behind_bg || !bg_opaque) {
            /* Bit 4 značí, že jde o pixel (relevantní
             * při finálním vykreslování na obrazovku) */
            *dest = px | 0x10;
        }
    }

    ++dest;
    tile >>= 4;
}
\end{minted}
\caption{Kód pro vykreslení spritu s řešením priority pořadí, Sprite 0 Hit a průhledností.}
\label{lst:render_sprites}
\end{listing}

Tato funkce řeší několik důležitých aspektu hardwaru NES:

\begin{enumerate}
    \item Sprity jsou prohledávány v opačném pořadí -- od nejnižší priority k nejvyšší.
    \item Je respektována průhlednost pixelů spritu -- pokud pixel spritu je průhledný, pixel v
        bufferu nebude ovlivněn.
    \item Pokud se vykresluje první v pořadí sprite, pixel pozadí není průhledný a přitom pixel
        spritu není 255. pixel, nastane podmínka Sprite 0 Hit a příslušný bit v \texttt{PPUSTAT}
        bude nastaven (který pak přečte emulovaný program a zjistí, zda Sprite 0 Hit skutečně
        nastal).
    \item Aby výsledný pixel spritu byl zapsán do bufferu, musí být vepředu (6. bit atributu
        spritu má být vypnut). Pokud není vepředu (odpovídající bit je zapnut), a přitom pixel
        pozadí není průhledný, pixel spritu bude zahozen.
\end{enumerate}

\subsection{Příprava dat pro vykreslování spritů}

Princip funkce \texttt{fetch\_spr\_tile()} pro přípravu řádku dlaždice spritu pro vykreslování je
stejný jako v případě řádku dlaždice pozadí. Vzhledem k tomu, že většina potřebné informace je již
zakódována v OAM, výsledný algoritmus je značně jednodušší. Tato funkce je také založena na balení
celého řádku dlaždice do 32bitového typu \texttt{TileRow} podobně jako \texttt{fetch\_spr\_tile()}, a
aktivně využívá LUT tabulky, zejména pro převracené řádky pixely, pokud je horizontální převracení
zapnuto.

Z důvodu zjednodušení se neřeší clipping. Toto však nemá zásadní vliv na základní funkčnost
emulátoru.

\subsection{Prohledávání OAM (Sprite Evaluation)}

Jak bylo zmíněno v podkapitole \ref{tp:sprite-eval}, systém NES dokáže na jednom řádku zobrazit
maximálně 8 spritů, a je nutné před samotným vykreslováním provést jejich výběr. Funkce
\texttt{queue\_sprites} (výpis \ref{lst:queue_sprites}) prochází paměť OAM a identifikuje objekty,
které protínají aktuálně emulovaný scanline.

\begin{listing}[htbp]
\begin{minted}{c}
for (unsigned i = 0; i < 64 && self->queued_sprites_count < 8; ++i) {
    uint8_t y = self->oam[i * 4];
    if (y >= 239) continue; /* Sprite mimo viditelnou oblast */

    unsigned row = self->scanline - y;
    if (row < h) { /* Sprite protíná aktuální řádek */
        Sprite *spr = &self->queued_sprites[self->queued_sprites_count];
        uint8_t tile = self->oam[i * 4 + 1];
        uint8_t attr = self->oam[i * 4 + 2];
        uint16_t pat_addr;

        spr->pos_y = y;
        spr->pos_x = self->oam[i * 4 + 3];
        spr->attr = attr;
        spr->is_0 = (i == 0); /* Důležité pro detekci Sprite 0 Hit */

        /* Výpočet jemné Y souřadnice s ohledem na vertikální převrácení */
        unsigned y_fine = (attr & 0x80) ? (h - 1 - row) : row;

        /* Výpočet adresy v Pattern Table */
        if (h == 8) { /* Režim 8x8 */
            uint16_t base = get_pattern_table_spr(self);
            pat_addr = base + (tile * 16) + y_fine;
        } else { /* Režim 8x16 (bit 0 určuje banku, zbytek index) */
            uint16_t base = (tile & 1) ? 0x1000 : 0x0000;
            uint8_t tile_idx = tile & 0xfe;
            if (y_fine >= 8) {
                ++tile_idx;
                y_fine -= 8;
            }
            pat_addr = base + (tile_idx * 16) + y_fine;
        }

        /* Přednačtení dat dlaždice pomocí LUT (stejně jako u pozadí) */
        spr->tile = fetch_spr_tile(self, pat_addr, attr);
        ++self->queued_sprites_count;
    }
\end{minted}
\caption{Implementace Sprite Evaluation. Funkce připravuje data pro 8 objektů na aktuálním řádku.}
\label{lst:queue_sprites}
\end{listing}

\newpage

Klíčové aspekty implementace:

\begin{itemize}
    \item Smyčka prochází všech 64 spritů (256 bytů) v pamětí OAM, ale ukončí se ihned po nalezení
        prvních osmi platných objektů. Tim je emulován limit původního hardwaru a zároveň šetřen
        procesorový čas mikrokontroléru.
    \item Tato implementace bere v povahu to, že systém NES měl sprity posunute o jeden pixel (viz
        \ref{tp:sprite-eval}).
    \item Funkce dynamicky přepíná mezi standardním režimem 8\texttimes8 a rozšířeným 8\texttimes16
        podle registru \texttt{PPU\_CTRL}. V režimu 8\texttimes16 funguje adresování
        odlišně -- nultý bit indexu dlaždice slouží jako volič banky (Pattern Table), zatímco
        zbylé bity určují konkrétní dvojici dlaždic.
    \item Pokud je v atributech nastaven příznak pro vertikální převracení (0x80), index řádku
        uvnitř dlaždice se invertuje. Tato operace probíhá před výpočtem adresy v CHR pamětí (pokud
        by převracení se aplikovalo při vykreslování, značně by to zkomplikovalo algoritmus a také
        výpočetní výkon).
\end{itemize}

